\chapter{Knee Pain}

\section*{Definition and Overview}
Knee pain is one of the most common reasons people visit their GP, physiotherapist, or emergency department. It can arise from injury, overuse, ageing, or disease affecting any of the structures around the joint, including bones, cartilage, ligaments, tendons, and bursae. The pain may be sharp or dull, sudden or gradual, and may be accompanied by swelling, stiffness, instability, or locking. Because the knee carries much of the body's weight and is essential for mobility, knee pain can significantly affect daily life. This chapter outlines the common causes of knee pain and how they are assessed and treated.

\section*{Epidemiology}
Knee pain affects people of all ages but the underlying causes differ across the lifespan. In children and adolescents, causes include patellofemoral pain, Osgood--Schlatter disease, and injury. In young and middle-aged adults, sports injuries to ligaments and menisci are common. In older adults, osteoarthritis is the leading cause of chronic knee pain, affecting a large proportion of people over 65, and the knee is the joint most commonly affected. Overweight and obesity, previous knee injury, and physically demanding occupations all increase the risk. Knee pain is a leading cause of disability worldwide.

\section*{Aetiology and Pathophysiology}
\begin{itemize}
    \item \textbf{Osteoarthritis:} progressive loss of joint cartilage with age, wear, obesity, and previous injury, leading to bone-on-bone contact, inflammation, and pain
    \item \textbf{Ligament injuries:} sprains or tears of the ACL, PCL, or collateral ligaments from trauma or twisting
    \item \textbf{Meniscal tears:} tears of the cartilage pads from twisting injuries or age-related degeneration
    \item \textbf{Patellofemoral pain:} pain around or behind the kneecap, common in runners and adolescents, related to tracking of the kneecap and muscle imbalance
    \item \textbf{Tendinopathy:} inflammation or degeneration of the patellar, quadriceps, or hamstring tendons from overuse
    \item \textbf{Bursitis:} inflammation of the fluid sacs around the knee, often from kneeling or overuse
    \item \textbf{Inflammatory arthritis:} rheumatoid arthritis, gout, or psoriatic arthritis affecting the knee
    \item \textbf{Infection:} septic arthritis, a medical emergency
    \item \textbf{Other causes:} bone tumours, referred pain from the hip or spine, and fractures
\end{itemize}

\section*{Clinical Features}
\begin{itemize}
    \item Pain that may be constant or only on movement, weight bearing, or stairs
    \item Stiffness, especially in the morning or after sitting (suggestive of arthritis) or after activity (suggestive of overuse)
    \item Swelling around the knee, which may be rapid (injury) or gradual (arthritis)
    \item Crepitus: a grating or cracking sensation with movement
    \item Giving way, instability, or a sensation that the knee will collapse
    \item Locking or catching, suggesting a meniscal tear or loose body
    \item Difficulty with walking, stairs, kneeling, or squatting
    \item Redness, warmth, and fever, which may indicate infection or inflammatory arthritis
\end{itemize}

\section*{Differential Diagnosis}
\begin{itemize}
    \item Osteoarthritis and other degenerative conditions
    \item Ligament, meniscal, and tendon injuries
    \item Patellofemoral pain syndrome
    \item Bursitis and tendinopathy
    \item Inflammatory arthritis: rheumatoid, gout, or psoriatic
    \item Septic arthritis (a medical emergency, especially with fever and a hot, swollen joint)
    \item Bone pathology: fracture, osteonecrosis, or rarely tumour
    \item Referred pain from the hip or lumbar spine
\end{itemize}

\section*{When to Seek Medical Care}
See a doctor if knee pain is severe, follows a significant injury, prevents weight bearing, is associated with locking or instability, or persists beyond a few days despite simple measures. Older adults with new, persistent knee pain and anyone with signs of inflammatory joint disease should also seek review.

\textbf{Call 000 (or local emergency services) for:}
\begin{itemize}
    \item An obviously deformed, dislocated, or angulated knee
    \item A hot, swollen, very painful knee with fever, which may be septic arthritis
    \item Inability to bear any weight or move the joint after trauma
    \item Numbness, coldness, or pallor of the lower leg or foot after an injury
    \item Sudden severe pain with significant swelling after a fall or impact
\end{itemize}

\section*{Diagnosis and Investigations}
\begin{itemize}
    \item \textbf{History:} the nature, onset, and location of pain; aggravating and relieving factors; mechanism of any injury; age; occupation; and activity level
    \item \textbf{Examination:} observation of alignment and swelling, palpation for tenderness, assessment of range of movement, and stability tests
    \item \textbf{X-ray:} for suspected fracture, or in older adults to assess arthritis
    \item \textbf{Magnetic resonance imaging (MRI):} for suspected ligament, meniscal, or cartilage damage when management depends on the diagnosis
    \item \textbf{Ultrasound:} for tendon, bursa, and superficial soft tissue problems
    \item \textbf{Blood tests:} inflammatory markers, uric acid, or autoantibody screens when inflammatory or crystal arthritis is suspected
    \item \textbf{Joint aspiration:} fluid is drawn from the knee to diagnose infection or gout
\end{itemize}

\section*{Management}
\begin{itemize}
    \item \textbf{Acute injuries:} rest, ice, compression, and elevation, with gradual resumption of activity
    \item \textbf{Physiotherapy:} strengthening, stretching, balance, and gait retraining benefit nearly all causes of knee pain
    \item \textbf{Pain relief:} paracetamol, topical anti-inflammatory gels, or short courses of oral anti-inflammatories
    \item \textbf{Osteoarthritis:} exercise, weight loss, activity modification, and, when needed, injections or surgery; joint replacement is considered when pain is severe and other measures have failed
    \item \textbf{Inflammatory arthritis:} disease-modifying medications under specialist care
    \item \textbf{Infection:} urgent antibiotics and joint drainage in hospital
    \item \textbf{Surgery:} arthroscopy for locking meniscal tears, ligament reconstruction, and, for advanced arthritis, partial or total knee replacement
    \item \textbf{Assistive devices:} supportive footwear, braces, or walking aids where appropriate
\end{itemize}

\section*{Complications}
\begin{itemize}
    \item Chronic pain and disability with reduced mobility and independence
    \item Loss of muscle strength and joint stiffness from inactivity
    \item Progression of osteoarthritis and joint deformity
    \item Falls, which are more likely in older people with painful, unstable knees
    \item Long-term absence from work or sport
    \item Rarely, missed infection or fracture leading to serious outcomes
\end{itemize}

\section*{Prognosis and Outcomes}
The outlook for knee pain depends on the cause. Overuse and minor injuries generally resolve with rest and rehabilitation within weeks. Osteoarthritis is a chronic condition, but exercise, weight management, and pacing can keep many people active for years; total knee replacement, when needed, is highly effective at relieving pain and restoring function. With an accurate diagnosis and appropriate treatment, most people can expect substantial improvement.

\section*{Prevention and Self-Care}
\begin{itemize}
    \item Keep the thigh and calf muscles strong through regular, low-impact exercise such as swimming, cycling, or walking
    \item Maintain a healthy weight to reduce load on the knee
    \item Warm up before sport and progress training loads gradually
    \item Wear supportive, appropriate footwear
    \item Use good technique and injury-prevention programmes in sport
    \item Avoid prolonged kneeling and high-impact exercise on painful knees
    \item Seek early treatment for injuries rather than "pushing through" pain
\end{itemize}

\section*{Sources and Further Reading}
The following reputable sources provide further information for healthcare students and patients seeking reliable, current advice about knee pain.
\begin{itemize}
    \item Healthdirect. \textit{Knee pain.} https://www.healthdirect.gov.au/knee-pain
    \item Arthritis Australia. \textit{Knee osteoarthritis information.} https://arthritisaustralia.com.au/
    \item Australian Physiotherapy Association. \textit{Knee pain resources.} https://australian.physio/
    \item NHS. \textit{Knee pain.} https://www.nhs.uk/conditions/knee-pain/
    \item Mayo Clinic. \textit{Knee pain.} https://www.mayoclinic.org/diseases-conditions/knee-pain/
    \item MSD Manual. \textit{Knee pain.} https://www.msdmanuals.com/
\end{itemize}
